\documentclass[tikz,border=4pt]{standalone}
\usepackage{ctex}
\usepackage{tikz}
\usetikzlibrary{arrows.meta, decorations.pathreplacing}
\begin{document}
\begin{tikzpicture}[scale=1.0]
% Number line
\draw[thick, -Stealth] (-5.5,0) -- (5.5,0) node[right] {$x$};

% Convergence interval
\draw[blue, very thick] (-3,0) -- (3,0);
\fill[blue!20] (-3,0) rectangle (3,0.15);

% Points
\fill[red] (-3,0) circle (3.5pt) node[below=4pt, font=\small] {$-R$};
\fill[green!60!black] (3,0) circle (3.5pt) node[below=4pt, font=\small] {$R$};
\fill (0,0) circle (2pt) node[below=4pt, font=\small] {$0$};

% Labels
\draw[decorate, decoration={brace,amplitude=6pt,mirror}, blue]
    (-3,-0.5) -- (3,-0.5) node[midway, below=8pt, font=\small, blue]
    {收敛区间 $(-R, R)$：绝对收敛};

% Left region
\draw[red!70, thick] (-5.5,0.3) -- (-3.2,0.3);
\node[red!70, font=\small] at (-4.5, 0.55) {$|x|>R$：发散};

% Right region
\draw[red!70, thick] (3.2,0.3) -- (5.5,0.3);
\node[red!70, font=\small] at (4.5, 0.55) {$|x|>R$：发散};

% Endpoint annotation for x=1 (example sum x^n/n)
\node[font=\small, align=center] at (-3, 1.2)
    {端点 $x=-R$\\需单独判};
\node[font=\small, align=center] at (3, 1.2)
    {端点 $x=R$\\需单独判};
\draw[-Stealth, gray] (-3,0.9) -- (-3,0.15);
\draw[-Stealth, gray] (3,0.9) -- (3,0.15);

% Example
\node[font=\small, gray] at (0, -1.2)
    {例：$\sum x^n/n$ → $R=1$，端点：$x=1$ 散，$x=-1$ 收 → 域 $[-1,1)$};
\end{tikzpicture}
\end{document}
